\section*{Computational Methods (full protocol)}

All total-energy and electronic-structure calculations use CP2K~\cite{cp2k2025} with PBE~\cite{Perdew1996generalized} and Grimme DFT-D3 (Becke--Johnson damping)~\cite{Grimme2011effect}.
Unless stated otherwise, biaxial strain scales the in-plane lattice vectors of the pristine equilibrium cell at fixed fractional coordinates (rigid-strain protocol), providing an upper bound on strain--doping coupling until ionic-relaxation benchmarks in Table~S3 converge.
Plane-wave cutoff is 300~Ry for $n\leq 5$ and 280~Ry for $n\geq 6$; $\mathcal{S}(n)$ comparisons at fixed $n$ use identical cutoffs.

Tetramer models ($4\times\mathrm{C}_{60}$, 240 sites) use B/N/P substitution at nominal 5\% (12--13 dopants from archived coordinates, seed~42).
Periodic scaling cells contain one dopant per C$_{60}$ ($n=1$--$8$).
Dopant sites are not re-optimized thermodynamically; fixed configurations probe mechanism rather than ensemble averages.

Grids: 24 tetramer strain points ($\epsilon\in\{-5,-2.5,0,2.5,5\}$\%); 12 tetramer electronic points; 40 periodic neutral-energy points at $\epsilon\in\{0,+3\}$\%; charged-state optimizations for Marcus schematics (Fig.~S5); IPR/$J$ factorial (Fig.~S6).
Linear $\alpha=\mathrm{d}E/\mathrm{d}\epsilon$ and substitution energies $E_f$ are defined in Table~S1.

\textbf{Non-additive coupling metric.}
We quantify strain--doping cross terms with the synergy order parameter $\mathcal{S}$:
\begin{equation}
\mathcal{S}(\epsilon,\delta)
= \Delta E(\epsilon,\delta) - \bigl[\Delta E(\epsilon,0) + \Delta E(0,\delta)\bigr],
\label{eq:synergy_order}
\end{equation}
where $\epsilon$ is biaxial strain, $\delta\in\{B,N,P\}$, and $\Delta E(\epsilon,\delta)$ is the energy per atom relative to the pristine, zero-strain reference at the same supercell size $n\times\mathrm{C}_{60}$.
Additive superposition gives $\mathcal{S}=0$; $\mathcal{S}<0$ indicates cooperative stabilization beyond the two-scan sum, whereas $\mathcal{S}>0$ indicates cooperative destabilization.
Exp.~10 values are computed with \texttt{c/simukit-sdc} (\texttt{sdc\_exp10\_results.json}; synergy audit table \texttt{sdc\_exp10\_synergy\_audit.json}).

\textbf{Localization and coupling (Exp.~4).}
Following Capobianco \emph{et al.}~\cite{Capobianco2024electron}, we report the inverse participation ratio (IPR) of the charged-state frontier orbital and the nearest-neighbor electronic coupling $J$ from tetramer dimer fragments.
IPR follows Electron SI Eq.~(S8), $N\sum_i|\psi_i|^4/(\sum_i|\psi_i|^2)^2$, on a Mulliken-projected orbital grid; $J$ is estimated from the ground--first-excited splitting of the anionic dimer at fixed interfullerene geometry (two-state model).
These observables are computed at the same PBE+D3 level as the energy scans and are \emph{not} directly comparable in absolute value to the Koopmans-compliant $J\simeq 75$--$81$~meV reported for pristine qHP networks~\cite{Capobianco2024electron}, but they provide an internally consistent trend under coupled $(\epsilon,\delta)$.

\textbf{Polaron binding (Exp.~9, in progress).}
Marcus reorganization energies follow the vertical--adiabatic decomposition~\cite{Capobianco2024electron}:
$\lambda^{\pm}=E_{\mathrm{vert}}^{\pm}-E_{\mathrm{opt}}^{\pm}$, with separate internal ($\lambda^{\pm}_{\mathrm{loc}}$) and environmental ($\lambda^{\pm}_{\mathrm{env}}$) contributions when geometry subsets are relaxed.
At the time of writing, 7/12 charged geometry optimizations have converged; vertical single-points at neutral geometry (needed for $\lambda$) are pending---values are not quoted in Results [live counts: \texttt{exp9\_polaron\_verification.json}; refresh via \texttt{post\_exp9\_converged.sh}].

